\documentclass{article}
% ---------------------------------------------------------------
% Póster de revisión de literatura (A0 vertical)
% Finanzas I (055D), UNCP, 2026-II. Compilar con pdfLaTeX.
% ---------------------------------------------------------------
\usepackage[paperwidth=841mm,paperheight=1189mm,margin=0mm]{geometry}
\usepackage[utf8]{inputenc}
\usepackage[T1]{fontenc}
\IfFileExists{spanish.ldf}{\usepackage[spanish,es-noshorthands,es-nodecimaldot]{babel}}{}
\usepackage{helvet}
\renewcommand{\familydefault}{\sfdefault}
\usepackage[letterspace=120]{microtype}
\usepackage{anyfontsize}
\usepackage{graphicx}
\usepackage{xcolor}
\usepackage[most]{tcolorbox}
\usepackage{enumitem}
\usepackage{multicol}
\usepackage{xurl}
\usepackage[hidelinks]{hyperref}
\usepackage{qrcode}
\urlstyle{same}
\pagestyle{empty}
\setlength{\parindent}{0pt}
\hyphenpenalty=10000 \exhyphenpenalty=10000

% ---------------- Paleta ----------------
\definecolor{navy}{HTML}{8E0C11}
\definecolor{navyLight}{HTML}{B8121A}
\definecolor{gold}{HTML}{00A3E0}
\definecolor{acento}{HTML}{A6E1F7}
\definecolor{goldSoft}{HTML}{E8F6FC}
\definecolor{teal}{HTML}{353435}
\definecolor{tealSoft}{HTML}{F1EEEE}
\definecolor{ink}{HTML}{231F20}
\definecolor{muted}{HTML}{5E5757}
\definecolor{line}{HTML}{E5D7D7}
\definecolor{page}{HTML}{F6F1F1}
\definecolor{mist}{HTML}{F7D3D5}
\pagecolor{page}

% ---------------- Tamaños (ajustados por póster) ----------------
\newcommand{\bodysize}{\fontsize{32}{42.2}\selectfont}
\newcommand{\cardtitlesize}{\fontsize{37.8}{45.4}\selectfont}
\newcommand{\smallsize}{\fontsize{26.2}{34.1}\selectfont}

% ---------------- Estilos de tarjeta ----------------
\tcbset{
  card/.style={enhanced, sharp corners, colback=white, colframe=line, boxrule=0.9mm,
    left=14mm, right=14mm, top=12mm, bottom=12mm, fontupper=\bodysize\color{ink}\raggedright,
    overlay={\fill[#1] (frame.north west) rectangle ([yshift=-5mm]frame.north east);}},
  card/.default=navy,
}
\newcommand{\cardtitle}[2][navy]{{\color{#1}\cardtitlesize\bfseries\textls[110]{\MakeUppercase{#2}}}\par\vspace{7mm}}
\newcommand{\bul}{\raisebox{0.18em}{\color{gold}\rule{0.5em}{0.5em}}}
\setlist[itemize]{label=\bul, leftmargin=1.25em, labelsep=0.6em, itemsep=0.55em, topsep=0pt, parsep=0pt}
\newcommand{\pill}[3]{\tcbox[on line, enhanced, arc=7mm, boxrule=0.8mm, colback=#1, colframe=#2,
  left=7mm, right=7mm, top=3mm, bottom=3mm, fontupper=\fontsize{26}{30}\selectfont\bfseries]{#3}}

\begin{document}

% ======================= ENCABEZADO =======================
\begin{tcolorbox}[enhanced, sharp corners, colback=navy, boxrule=0pt, width=\paperwidth,
  left=34mm, right=34mm, top=20mm, bottom=20mm,
  overlay={\fill[gold] (frame.south west) rectangle ([yshift=5mm]frame.south east);
           \fill[navyLight] ([xshift=-160mm]frame.north east) -- (frame.north east) -- ([yshift=-160mm]frame.north east) -- cycle;}]
\begin{minipage}[c]{92mm}\centering\includegraphics[width=\linewidth]{Escudo_UNCP.png}\end{minipage}\hfill
\begin{minipage}[c]{555mm}\centering
{\color{acento}\fontsize{25}{30}\selectfont\bfseries\textls[160]{REVISIÓN DE LITERATURA \quad|\quad FINANZAS I \quad|\quad TEMA 61}\par}
\vspace{10mm}
{\color{white}\fontsize{66}{74.6}\selectfont\bfseries Algorithmic and high-frequency trading in Borsa Istanbul\par}
\vspace{8mm}
{\color{mist}\fontsize{38}{46.4}\selectfont\itshape Negociación algorítmica y de alta frecuencia en Borsa Istanbul\par}
\vspace{10mm}
{\color{white}\fontsize{29}{35}\selectfont\bfseries Oguz Ersan y Cumhur Ekinci\par}
\vspace{9mm}
\pill{white}{white}{\color{navy}Q1}\hspace{6mm}\pill{navy}{white}{\color{white}\mdseries \textit{Borsa Istanbul Review}, 16(4), 233-248 (2016)}
\end{minipage}\hfill
\begin{minipage}[c]{92mm}\centering\includegraphics[width=\linewidth]{Logo_FE.png}\end{minipage}
\end{tcolorbox}
\vspace{-1mm}
\begin{tcolorbox}[enhanced, sharp corners, colback=goldSoft, boxrule=0pt, width=\paperwidth,
  left=34mm, right=34mm, top=8mm, bottom=8mm, halign=center]
{\color{navy}\fontsize{27}{36}\selectfont \textbf{Póster elaborado por:} Meglinho Sneyder Rojas Poma \enspace\textbar\enspace Código 2024200522B\par \fontsize{24}{32}\selectfont Escuela Profesional de Economía, Universidad Nacional del Centro del Perú \enspace\textbar\enspace Finanzas I (055D), 2026-II\par}
\end{tcolorbox}

\vfill
% ======================= CUERPO =======================
\hspace*{34mm}\begin{minipage}{773mm}
\begin{tcolorbox}[enhanced, sharp corners, colback=navy, boxrule=0pt, left=16mm, right=16mm, top=12mm, bottom=11mm,
  fontupper=\bodysize\raggedright, overlay={\fill[gold] (frame.north west) rectangle ([yshift=-5mm]frame.north east);}]
\begin{minipage}[t]{174.8mm}\raggedright{\color{acento}\smallsize\bfseries\textls[110]{\MakeUppercase{Mercado}}}\par\vspace{2mm}{\color{white}Borsa Istanbul (Turquía)}\end{minipage}\hfill\begin{minipage}[t]{174.8mm}\raggedright{\color{acento}\smallsize\bfseries\textls[110]{\MakeUppercase{Muestra}}}\par\vspace{2mm}{\color{white}354 días de negociación}\end{minipage}\hfill\begin{minipage}[t]{174.8mm}\raggedright{\color{acento}\smallsize\bfseries\textls[110]{\MakeUppercase{Periodo}}}\par\vspace{2mm}{\color{white}Enero de 2013 a mayo de 2014}\end{minipage}\hfill\begin{minipage}[t]{174.8mm}\raggedright{\color{acento}\smallsize\bfseries\textls[110]{\MakeUppercase{Método}}}\par\vspace{2mm}{\color{white}Medidas de mensajes y órdenes}\end{minipage}
\end{tcolorbox}
\end{minipage}

\vfill
\hspace*{34mm}\begin{minipage}{773mm}
\begin{tcbraster}[raster columns=3, raster equal height, raster column skip=16mm]
\begin{tcolorbox}[card]\cardtitle{Contexto}La negociación algorítmica, hecha por programas de computadora, crece con la tecnología. La de alta frecuencia es una parte de ella que envía muchas órdenes en milisegundos.\end{tcolorbox}
\begin{tcolorbox}[card]\cardtitle{Problema}No se conocía cuánta negociación algorítmica y de alta frecuencia había en un mercado emergente como Borsa Istanbul.\end{tcolorbox}
\begin{tcolorbox}[card]\cardtitle{Objetivo}Medir el nivel de negociación algorítmica y de alta frecuencia en Borsa Istanbul y proponer una medida ampliada de la alta frecuencia.\end{tcolorbox}
\end{tcbraster}
\end{minipage}

\vfill
\hspace*{34mm}\begin{minipage}{773mm}
\begin{tcolorbox}[card]\cardtitle{Metodología}
\begin{minipage}[t]{0.48\linewidth}
\begin{itemize}
\item La negociación algorítmica se mide con los mensajes por minuto y la medida de Hendershott y otros.
\item La alta frecuencia se detecta con mensajes enlazados.
\item Se propone una medida ampliada que capta varias estrategias.
\end{itemize}
\end{minipage}\hfill
\begin{minipage}[t]{0.48\linewidth}
\begin{itemize}
\item Se analizan 85 millones de órdenes y 243 millones de mensajes de las acciones del BIST 100.
\item Se estudian tipos de mensajes, secuencias de órdenes y formas de cancelación o ejecución.
\end{itemize}
\end{minipage}
\end{tcolorbox}
\end{minipage}

\vfill
\hspace*{34mm}\begin{minipage}{773mm}
\begin{tcolorbox}[card]\cardtitle{Figura clave del estudio}
\begin{minipage}[c]{330.7mm}\centering
\fcolorbox{line}{white}{\includegraphics[width=\dimexpr330.7mm-2\fboxsep-2\fboxrule\relax]{figura.png}}
\end{minipage}\hfill
\begin{minipage}[c]{\dimexpr\linewidth-330.7mm-18mm\relax}\raggedright
{\color{navy}\bfseries Figura 4.}\enspace{\itshape Características de las acciones según su nivel de alta frecuencia (HFT): (a) capitalización, (b) liquidez y (c) volatilidad. La línea roja es el nivel de HFT y la negra la característica, para 120 acciones.}\par\vspace{8mm}
{\color{teal}\smallsize\bfseries\textls[110]{CÓMO LEERLA}}\par\vspace{4mm}
\begin{itemize}
\item La mayoría de acciones tiene niveles de HFT repartidos de forma casi uniforme entre 2 \% y 7 \%.
\item El decil con HFT excesivo, de 8 \% a 14 \% en promedio, son casi todas acciones de empresas pequeñas.
\item Esas acciones también son relativamente ilíquidas.
\end{itemize}
\vspace{8mm}
{\color{muted}\smallsize Figura tomada de Ersan y Ekinci (2016).}
\end{minipage}
\end{tcolorbox}
\end{minipage}

\vfill
\hspace*{34mm}\begin{minipage}{773mm}
\begin{tcolorbox}[card=teal]\cardtitle[teal]{Resultados}
\begin{tcolorbox}[enhanced, sharp corners, colback=tealSoft, boxrule=0pt, left=14mm, right=14mm, top=10mm, bottom=11mm,
  borderline west={4mm}{0mm}{teal}, fontupper=\bodysize\color{ink}\raggedright]
{\color{teal}\smallsize\bfseries\textls[110]{HALLAZGO CENTRAL}}\par\vspace{4mm}
{\bfseries La negociación algorítmica crece y la alta frecuencia representa alrededor del 6 \% de las órdenes.}
\end{tcolorbox}
\vspace{11mm}
\begin{minipage}[t]{0.48\linewidth}
\begin{itemize}
\item El nivel en el BIST 100 es parecido al de las grandes empresas del NASDAQ entre 2003 y 2005.
\item La alta frecuencia es mayor en las órdenes grandes (11.96 \%).
\end{itemize}
\end{minipage}\hfill
\begin{minipage}[t]{0.48\linewidth}
\begin{itemize}
\item También es mayor en las órdenes de gestoras de fondos (10.40 \%) y tras mejorar la plataforma y reducir el tick.
\item La medida algo\_trad promedio del BIST 100 está entre 18 y 13, con tendencia al alza de la negociación algorítmica.
\end{itemize}
\end{minipage}
\end{tcolorbox}
\end{minipage}

\vfill
\hspace*{34mm}\begin{minipage}{773mm}
\begin{tcbraster}[raster columns=3, raster equal height, raster column skip=16mm]
\begin{tcolorbox}[card, raster multicolumn=2]\cardtitle{Conclusiones}La tecnología y las reglas de la bolsa determinan cuánto se negocia de forma automatizada, lo que influye en la liquidez del mercado.\end{tcolorbox}
\begin{tcolorbox}[card=gold, colback=goldSoft]\cardtitle{Aporte al tema 61}Sirve de contraste con la BVL: un mercado con mucha negociación automatizada tiene precios que cambian todo el tiempo, mientras que en Lima muchos precios se repiten por la baja negociación.\end{tcolorbox}
\end{tcbraster}
\end{minipage}

\vfill
% ======================= PIE =======================
\begin{tcolorbox}[enhanced, sharp corners, colback=navy, boxrule=0pt, width=\paperwidth,
  left=34mm, right=34mm, top=14mm, bottom=11mm,
  overlay={\fill[gold] (frame.north west) rectangle ([yshift=-5mm]frame.north east);}]
\begin{minipage}[c]{245mm}\raggedright
{\color{acento}\smallsize\bfseries\textls[110]{FUENTE DE DATOS DEL ESTUDIO}}\par\vspace{5mm}
{\color{white}\smallsize Mensajes de órdenes de Borsa Istanbul durante 354 días de negociación, de 2013 a 2014.}
\end{minipage}\hfill
\begin{minipage}[c]{385mm}\raggedright
{\color{acento}\smallsize\bfseries\textls[110]{REFERENCIA}}\par\vspace{5mm}
{\color{white}\smallsize\hangindent=1.2em\hangafter=1 Ersan, O., y Ekinci, C. (2016). Algorithmic and high-frequency trading in Borsa Istanbul. \textit{Borsa Istanbul Review}, \textit{16}(4), 233-248. \url{https://doi.org/10.1016/j.bir.2016.09.005}\par}
\end{minipage}\hfill
\begin{minipage}[c]{100mm}\centering
\tcbox[enhanced, sharp corners, colback=white, boxrule=0pt, left=3mm, right=3mm, top=3mm, bottom=3mm]{\qrcode[hyperlink, height=46mm]{https://doi.org/10.1016/j.bir.2016.09.005}}\par\vspace{4mm}
{\color{acento}\fontsize{17}{21}\selectfont\bfseries\textls[100]{ESCANEA Y LEE EL ARTÍCULO}}
\end{minipage}
\end{tcolorbox}
\end{document}
